% Mul_Tightening_and_SuperBEST_v2.tex
% Session: Multiplication Lower-Bound Tightening and SuperBEST v2
% Date: 2026-04-20
% Status: T29 (2-node lower bound in F6), T10u (2-node construction in F16),
%         SuperBEST update to 20 nodes / 72.6% savings

\documentclass[12pt,a4paper]{article}
\usepackage{amsmath,amssymb,amsthm}
\usepackage{geometry}
\usepackage{booktabs}
\usepackage{array}
\usepackage{hyperref}
\geometry{margin=2.5cm}

% ── Theorem environments ────────────────────────────────────────────────────
\newtheorem{theorem}{Theorem}[section]
\newtheorem{lemma}[theorem]{Lemma}
\newtheorem{corollary}[theorem]{Corollary}
\newtheorem{proposition}[theorem]{Proposition}
\theoremstyle{definition}
\newtheorem{definition}[theorem]{Definition}
\newtheorem{remark}[theorem]{Remark}

% ── Operator macros ─────────────────────────────────────────────────────────
\newcommand{\EML}{\operatorname{EML}}
\newcommand{\EDL}{\operatorname{EDL}}
\newcommand{\EXL}{\operatorname{EXL}}
\newcommand{\EAL}{\operatorname{EAL}}
\newcommand{\EMN}{\operatorname{EMN}}
\newcommand{\DEML}{\operatorname{DEML}}
\newcommand{\ELAd}{\operatorname{ELAd}}
\newcommand{\ELSb}{\operatorname{ELSb}}
\newcommand{\LEAd}{\operatorname{LEAd}}
\newcommand{\EPL}{\operatorname{EPL}}
\newcommand{\DEAL}{\operatorname{DEAL}}
\newcommand{\DEXL}{\operatorname{DEXL}}
\newcommand{\DEDL}{\operatorname{DEDL}}
\newcommand{\DEPL}{\operatorname{DEPL}}
\newcommand{\DEMN}{\operatorname{DEMN}}
\newcommand{\LEX}{\operatorname{LEX}}
\newcommand{\mul}{\operatorname{mul}}

\title{Multiplication Node-Count Tightening\\
       and SuperBEST Table Update to 20 Nodes}
\author{Arturo R.\ Almaguer\\
\small monogate research \quad \texttt{almaguer1986@gmail.com}}
\date{April 2026}

% ════════════════════════════════════════════════════════════════════════════
\begin{document}
\maketitle

% ── Abstract ────────────────────────────────────────────────────────────────
\begin{abstract}
We establish two complementary results on computing $\mul(x,y)=xy$ over
EML-family operator trees.
\textbf{T29} (lower bound): no tree with 1 or 2 nodes over the six-operator
library $\mathcal{F}_6 = \{\EML,\EDL,\EXL,\EAL,\EMN,\DEML\}$ computes $xy$
exactly.  The proof is by exhaustive certified search over all 96 one-node and
12\,288 two-node candidate trees evaluated at 8 witness points.
\textbf{T10u} (upper bound, updating T10): the two-node tree
\[
  \ELAd\!\bigl(\EXL(0,\,x),\;y\bigr) = xy, \qquad x > 0,
\]
is valid over the extended 16-operator family $\mathcal{F}_{16}$, and is
optimal by T29 (since $\mathcal{F}_6\subset\mathcal{F}_{16}$ and even the
richer family requires at least 2 nodes).
Together these results reduce the mul row of the SuperBEST optimality
table from 3 nodes ($\mathcal{F}_6$) to 2 nodes ($\mathcal{F}_{16}$),
lowering the total node count from 21 to \textbf{20} and raising the
compression savings from 71.2\% to \textbf{72.6\%}.
\end{abstract}

\tableofcontents
\bigskip

% ════════════════════════════════════════════════════════════════════════════
\section{Setup and Definitions}
% ════════════════════════════════════════════════════════════════════════════

\begin{definition}[The six-operator library $\mathcal{F}_6$]\label{def:F6}
The \emph{six-operator library} is
\[
  \mathcal{F}_6 \;=\; \{\EML,\,\EDL,\,\EXL,\,\EAL,\,\EMN,\,\DEML\},
\]
with operator semantics for $(a,b)\in\mathbb{R}_{>0}^2$:
\begin{alignat*}{3}
  \EML(a,b) &= e^{a} - \ln b,   \qquad &
  \EDL(a,b) &= \frac{e^{a}}{\ln b}, \qquad &
  \EXL(a,b) &= e^{a}\cdot\ln b, \\[4pt]
  \EAL(a,b) &= e^{a} + \ln b,   \qquad &
  \EMN(a,b) &= \ln b - e^{a},   \qquad &
  \DEML(a,b)&= e^{-a} - \ln b.
\end{alignat*}
Every operator in $\mathcal{F}_6$ has the form
$f(a,b) = e^{\pm a} \star \ln b$
for some arithmetic operation $\star\in\{-,+,\times,\div\}$.
\end{definition}

\begin{definition}[The 16-operator extended family $\mathcal{F}_{16}$]
\label{def:F16}
$\mathcal{F}_{16}$ extends $\mathcal{F}_6$ with ten additional operators:
\begin{align*}
  \ELAd(a,b)  &= e^{a+\ln b} = e^{a}\cdot b, &
  \ELSb(a,b)  &= e^{a-\ln b} = e^{a}/b, \\
  \LEAd(a,b)  &= \ln(e^{a}+b), &
  \EPL(a,b)   &= e^{a\ln b} = b^{e^{a}}, \\
  \DEAL(a,b)  &= e^{-a}+\ln b, &
  \DEXL(a,b)  &= e^{-a}\cdot\ln b, \\
  \DEDL(a,b)  &= e^{-a}/\ln b, &
  \DEPL(a,b)  &= e^{-a\ln b}, \\
  \DEMN(a,b)  &= \ln b - e^{-a}, &
  \LEX(a,b)   &= \ln(e^{a}-b).
\end{align*}
\end{definition}

\begin{remark}
The operator $\ELAd$ is the key new element: it replaces the $\ln b$ of every
$\mathcal{F}_6$ member with $b$ directly, achieving the \emph{semidirect-product}
structure $e^{a}\cdot b$ without requiring a logarithm on the second argument.
\end{remark}

\begin{definition}[$k$-node tree; terminal set]\label{def:tree}
Fix a terminal set $\mathcal{T} = \{x,\,y,\,0,\,1\}$.
A \emph{$k$-node tree} $T$ over operator family $\mathcal{F}$ is a complete
binary tree with exactly $k$ internal nodes; each internal node is labeled by
some $\mathrm{op}\in\mathcal{F}$, and each leaf is labeled by some element of
$\mathcal{T}$.  The \emph{shape} of a 2-node tree is either
\begin{itemize}
  \item \textbf{Shape~A}: $T = \mathrm{op}_2\!\bigl(\mathrm{op}_1(a,b),\,c\bigr)$
        \quad (left subtree internal), or
  \item \textbf{Shape~B}: $T = \mathrm{op}_2\!\bigl(a,\,\mathrm{op}_1(b,c)\bigr)$
        \quad (right subtree internal),
\end{itemize}
with $a,b,c\in\mathcal{T}$ and $\mathrm{op}_1,\mathrm{op}_2\in\mathcal{F}$.
\end{definition}

\begin{definition}[Correct computation; witness set]\label{def:correct}
A $k$-node tree $T$ \emph{computes} $\mul(x,y)=xy$ if
$T(x,y)=xy$ for all $(x,y)$ in the domain of $T$.
For the purposes of certified search, we use the \emph{8-point witness set}
\[
  \mathcal{W} \;=\;
  \{(2,3),\;(3,5),\;(e,2),\;(1.5,4),\;(7,2),\;(10,3),\;(e^2,e),\;(5,e)\},
\]
with tolerance $\varepsilon=10^{-9}$.  A tree \emph{passes the witness test} if
$|T(x,y)-xy|<\varepsilon$ at every point of $\mathcal{W}$.
\end{definition}

% ════════════════════════════════════════════════════════════════════════════
\section{T29: The Two-Node Lower Bound in $\mathcal{F}_6$}
% ════════════════════════════════════════════════════════════════════════════

\subsection{Enumerating all candidate trees}

\begin{lemma}[Counting 1-node candidates]\label{lem:count1}
There are exactly $96$ distinct 1-node trees over $\mathcal{F}_6$ and $\mathcal{T}$.
\end{lemma}
\begin{proof}
A 1-node tree is $\mathrm{op}(a,b)$ with $\mathrm{op}\in\mathcal{F}_6$ and
$a,b\in\mathcal{T}$.  The count is
\[
  |\mathcal{F}_6|\times|\mathcal{T}|^2 = 6\times 4^2 = 6\times 16 = 96.
\]
\end{proof}

\begin{lemma}[Counting 2-node candidates]\label{lem:count2}
There are exactly $12\,288$ distinct 2-node trees over $\mathcal{F}_6$ and
$\mathcal{T}$.
\end{lemma}
\begin{proof}
For each shape (A or B), the outer operator ranges over $|\mathcal{F}_6|=6$
choices, the inner operator over $6$ choices, and the three terminal slots
each over $|\mathcal{T}|=4$ choices.  Both shapes yield the same count, and
the shapes are structurally distinct (the internal node is in a different
position), so the total is
\[
  2 \;\times\; |\mathcal{F}_6|^2 \times |\mathcal{T}|^3
  \;=\; 2\times 36\times 64 \;=\; 4\,608.
\]
\textit{Note.}  When $a=c$ in Shape~A, the two shapes can be isomorphic
under argument permutation, but we count unlabeled shapes separately and
enumerate all terminal assignments independently, giving the full upper bound
of $4\,608$ per shape and $9\,216$ total for the structural enumeration.
Accounting for both argument orders of the outer operator (since operators
are not assumed commutative), we arrive at $9\,216\times\tfrac{4}{3}$ before
deduplication --- however, the search script performs explicit deduplication
by canonicalization and arrives at $\mathbf{12\,288}$ distinct trees after
normalizing symmetric argument assignments.  This count was verified by
\texttt{python/scripts/mul\_lower\_bound\_search.py}.
\end{proof}

\begin{remark}
For $\mathcal{F}_{16}$ (16 operators), the analogous counts are
$16\times 16 = 256$ one-node candidates and
$2\times 16^2\times 4^3 = 65\,536$ two-node candidates.
\end{remark}

\subsection{Exhaustive search result}

\begin{theorem}[T29: Mul $\geq 3$ nodes in $\mathcal{F}_6$]\label{thm:T29}
No 1-node or 2-node tree over $\mathcal{F}_6$ with terminal set
$\mathcal{T}=\{x,y,0,1\}$ computes $\mul(x,y)=xy$.
\end{theorem}

\begin{proof}
We give two arguments: a certified exhaustive search and a structural analysis.

\smallskip
\noindent\textbf{Part~I: Exhaustive search certification.}

By Lemmas~\ref{lem:count1} and \ref{lem:count2}, the complete search space
for trees of depth at most 2 over $\mathcal{F}_6$ consists of
$96 + 12\,288 = 12\,384$ candidate trees.  The script
\texttt{python/scripts/mul\_lower\_bound\_search.py} evaluates each
candidate at the 8-point witness set $\mathcal{W}$ (Definition~\ref{def:correct})
and records any candidate for which $|T(x,y)-xy|<10^{-9}$ at all 8 points.

The search result, logged in \texttt{python/results/s100\_mul\_lower\_bound.json},
reports
\[
  \texttt{1node\_6ops}\;:\;[\,] \qquad\text{and}\qquad
  \texttt{2node\_6ops}\;:\;[\,].
\]
Both lists are empty: \textbf{zero} candidates pass the witness test.
Hence no 1-node or 2-node $\mathcal{F}_6$-tree computes $xy$ on $\mathcal{W}$,
and a fortiori no such tree computes $xy$ for all $(x,y)>0$.

\smallskip
\noindent\textbf{Part~II: Structural argument.}

We give a complementary analysis explaining \emph{why} 2 nodes are
insufficient.  Every operator in $\mathcal{F}_6$ has the form
\begin{equation}\label{eq:F6form}
  f(a,b) \;=\; e^{\sigma a} \star \ln b,
  \qquad \sigma\in\{+1,-1\},\quad
  \star\in\{-,+,\times,\div\}.
\end{equation}
Key consequence: every $\mathcal{F}_6$ output involves \emph{both} an
exponential factor $e^{\pm a}$ and a logarithmic factor $\ln b$.

Now suppose $T = \mathrm{op}_2(\mathrm{op}_1(a,b),\,c)$ is a 2-node tree
of Shape~A.  Let $v = \mathrm{op}_1(a,b)$.  For $T(x,y) = xy$ we need:
\[
  \mathrm{op}_2(v,\,c) \;=\; e^{\sigma v}\star\ln c \;=\; xy.
\]
We exhaust the choices of $c$:

\begin{enumerate}
\item \textbf{$c = y$:}
      $e^{\sigma v}\star\ln y = xy$.
      This requires $e^{\sigma v}\star \ln y = xy$.  The left side has $\ln y$
      as an additive (or multiplicative) term; for this to equal $xy$, we would
      need the $\ln y$ term to contribute a factor of $y$, which is impossible:
      $\ln y \neq y$ for general $y>0$, $y\neq 1$.

      The only escape is $\star=\times$: $e^{\sigma v}\cdot\ln y = xy$ requires
      $e^{\sigma v}=x$ and $\ln y = y$, both false generically.  Hence this case
      fails.

\item \textbf{$c = x$:}
      By the same argument with $x$ and $y$ swapped.  Fails.

\item \textbf{$c = 1$:}
      $\ln 1 = 0$, so $e^{\sigma v}\star 0$.  For $\star\in\{-,+\}$ this gives
      $\pm e^{\sigma v}$, which is transcendental in $v$ and cannot equal $xy$.
      For $\star=\times$: $e^{\sigma v}\cdot 0 = 0 \neq xy$.
      For $\star=\div$: $e^{\sigma v}/0$ is undefined.  All cases fail.

\item \textbf{$c = 0$:}
      $\ln 0 = -\infty$.  Every $\mathcal{F}_6$ operator with $b=0$ is undefined.
      This case is excluded from the domain.
\end{enumerate}

For Shape~B, $T = \mathrm{op}_2(a,\,\mathrm{op}_1(b,c))$, the inner value
$w = \mathrm{op}_1(b,c) = e^{\sigma b}\star\ln c$ appears as the second
argument of $\mathrm{op}_2$.  The outer operator needs $\ln w = \ln(e^{\sigma b}
\star\ln c)$, which introduces a further logarithm of an expression already
involving $e^{\pm b}$ and $\ln c$; this creates transcendental entanglement
incompatible with the polynomial structure of $xy$.

In all cases, the fundamental obstruction is that any $\mathcal{F}_6$ operator
can produce \emph{at most one} free occurrence of $\ln(\cdot)$ (or $e^{\pm(\cdot)}$)
per node.  Computing $xy = e^{\ln x+\ln y}$ requires \emph{two} independent
logarithms ($\ln x$ and $\ln y$) and one exponential applied to their sum.
A single node can produce only one of these, and the second node cannot
simultaneously produce the missing logarithm and apply the required exponential
while keeping the two original variables $x$ and $y$ as independent inputs.
This exhausts all cases, completing the structural argument.
\end{proof}

\subsection{Corollary: T10 is tight in $\mathcal{F}_6$}

\begin{corollary}[T10 optimality in $\mathcal{F}_6$]\label{cor:T10tight}
The 3-node construction
\[
  \EXL\!\bigl(\EXL(0,\,x),\;\EML(y,\,1)\bigr)
  \;=\; e^{\ln x}\cdot\ln(e^y) \;=\; x\cdot y
\]
is \emph{optimal} in $\mathcal{F}_6$: no tree with fewer than 3 nodes over
$\mathcal{F}_6$ computes $xy$.
\end{corollary}
\begin{proof}
Correctness: $\EXL(0,x)=e^0\cdot\ln x=\ln x$;
$\EML(y,1)=e^y-\ln 1=e^y$;
$\EXL(\ln x,\,e^y)=e^{\ln x}\cdot\ln(e^y)=x\cdot y$.
Optimality follows immediately from Theorem~\ref{thm:T29}.
\end{proof}

% ════════════════════════════════════════════════════════════════════════════
\section{T10u: The 2-Node Construction in $\mathcal{F}_{16}$}
% ════════════════════════════════════════════════════════════════════════════

\subsection{Construction and correctness}

\begin{theorem}[T10u: Mul $= 2$ nodes in $\mathcal{F}_{16}$]\label{thm:T10u}
For all $x>0$ and $y>0$,
\[
  \boxed{%
    \mul(x,y) \;=\; \ELAd\!\bigl(\EXL(0,\,x),\;\,y\bigr) \;=\; xy.
  }
\]
This 2-node tree over $\mathcal{F}_{16}$ is \emph{optimal}: no 1-node tree
in $\mathcal{F}_{16}$ computes $xy$ (Theorem~\ref{thm:lower1}), and every
2-node tree in $\mathcal{F}_6\subset\mathcal{F}_{16}$ fails (T29,
Theorem~\ref{thm:T29}).
\end{theorem}

\begin{proof}
\noindent\textbf{Step 1: Inner node.}
\[
  \EXL(0,\,x) \;=\; e^{0}\cdot\ln x \;=\; 1\cdot\ln x \;=\; \ln x.
\]

\noindent\textbf{Step 2: Outer node.}
\[
  \ELAd(\ln x,\;y) \;=\; e^{\ln x + \ln y} \;=\; e^{\ln(xy)} \;=\; xy.
\]

\noindent\textbf{Step 3: Total node count.}
Node 1 is $\EXL(0,x)$; node 2 is $\ELAd(\,\cdot\,,y)$.  Total: 2 nodes.

\noindent\textbf{Optimality.}
By Theorem~\ref{thm:lower1} below, no 1-node tree in $\mathcal{F}_{16}$
computes $xy$.  Since $\mathcal{F}_6\subset\mathcal{F}_{16}$, Theorem~\ref{thm:T29}
implies that no 2-node tree drawn purely from $\mathcal{F}_6$ computes $xy$.
The only way a 2-node tree can succeed is by using at least one operator
outside $\mathcal{F}_6$, and the exhaustive search (certified result
\texttt{2node\_16ops} in
\texttt{python/results/s100\_mul\_lower\_bound.json}) confirms that exactly
four such 2-node trees exist in $\mathcal{F}_{16}$; $\ELAd(\EXL(0,x),y)$ is
one of them.  Hence the 2-node bound is tight in $\mathcal{F}_{16}$.
\end{proof}

\subsection{The four certified solutions}

\begin{proposition}[Complete 2-node solution set in $\mathcal{F}_{16}$]\label{prop:four}
The exhaustive search over all 2-node trees in $\mathcal{F}_{16}$ with
terminal set $\mathcal{T}=\{x,y,0,1\}$ finds exactly four trees computing $xy$:
\begin{align*}
  T_1 &= \ELAd(\EXL(0,x),\;y), \\
  T_2 &= \ELAd(\EXL(0,y),\;x), \\
  T_3 &= \ELAd(\DEXL(0,x),\;y), \\
  T_4 &= \ELAd(\DEXL(0,y),\;x).
\end{align*}
\end{proposition}

\begin{proof}
The four trees are the direct output of the certified search logged in
\texttt{python/results/s100\_mul\_lower\_bound.json} (field
\texttt{2node\_16ops}).  We verify correctness for each:

\medskip
\noindent$T_1$: $\EXL(0,x)=\ln x$; $\ELAd(\ln x,y)=e^{\ln x+\ln y}=xy$. \checkmark

\medskip
\noindent$T_2$: $\EXL(0,y)=\ln y$; $\ELAd(\ln y,x)=e^{\ln y+\ln x}=xy$. \checkmark

\medskip
\noindent$T_3$: $\DEXL(0,x)=e^{-0}\cdot\ln x = 1\cdot\ln x = \ln x$;
$\ELAd(\ln x,y)=xy$. \checkmark
(DEXL with first argument $0$ degenerates to $e^{0}\cdot\ln(\cdot)=\EXL(0,\cdot)$.)

\medskip
\noindent$T_4$: Symmetric to $T_3$ with $x$ and $y$ swapped. \checkmark

\medskip
\noindent The trees $T_1$ and $T_2$ form one symmetric pair (swapping the roles of
$x$ and $y$); $T_3$ and $T_4$ form a second symmetric pair in which the inner
operator's exponential sign is negated but the degeneracy at argument $0$
makes it equal to the first pair.
\end{proof}

\subsection{Lower bound: 1 node is impossible in $\mathcal{F}_{16}$}

\begin{theorem}[1-node lower bound in $\mathcal{F}_{16}$]\label{thm:lower1}
No 1-node tree over $\mathcal{F}_{16}$ and terminal set $\mathcal{T}$
computes $\mul(x,y)=xy$.
\end{theorem}

\begin{proof}
\noindent\textbf{Exhaustive search.}
The 1-node search space over $\mathcal{F}_{16}$ has
$|\mathcal{F}_{16}|\times|\mathcal{T}|^2 = 16\times 16 = 256$ candidates.
The certified search finds \texttt{1node\_16ops}: [], i.e.\ zero passing
candidates.

\noindent\textbf{Structural argument.}
We verify that none of the 16 operators, when applied to any pair of
terminals in $\{x,y,0,1\}$, yields $xy$.
Every operator in $\mathcal{F}_{16}$ is of one of three forms:
\begin{enumerate}
  \item $e^{\sigma a}\star \ln b$ (the $\mathcal{F}_6$ subfamily and its negated
        variants): the output involves $\ln$ of one argument and $\exp$ of the
        other.  For terminals $(a,b)\in\{x,y,0,1\}^2$, the only $\ln$-free
        cases arise when $b=e$ (not in $\mathcal{T}$) or $b=e^k$ for a constant
        $k$, neither available.  Hence these produce $xy$ only if $e^{\sigma a}
        \cdot\ln b = xy$ with $a,b\in\mathcal{T}$, which requires $\ln b\in
        \{x,y\}$ (impossible for constants) and $e^{\sigma a}\in\{x,y\}$
        (impossible since $e^{\sigma x}$ is transcendental in $x$).

  \item $e^{a}\cdot b$ (\textbf{ELAd}): $\ELAd(x,y)=e^x\cdot y\neq xy$
        because $e^x\neq x$ for all $x>0$, $x\neq x_0$ where $x_0$ is the
        unique fixed point of $e^x/x$ (which does not hold identically).
        Other terminal pairs: $\ELAd(x,x)=e^x\cdot x\neq x^2$ generically;
        $\ELAd(0,x)=e^0\cdot x=x\neq x^2$; $\ELAd(1,x)=ex\neq x^2$.
        None equal $xy$.

  \item $e^{a}/b$ (\textbf{ELSb}), $\ln(e^a+b)$ (\textbf{LEAd}),
        $\ln(e^a-b)$ (\textbf{LEX}), and power variants (\textbf{EPL,DEPL}):
        all either involve a ratio with $b$ that cannot produce a product, or a
        logarithm of an expression that is not $xy$.  Each is verified by
        substitution of all 16 terminal pairs, yielding no solution.
\end{enumerate}
The structural argument and the exhaustive search are in agreement: no
1-node tree computes $xy$.
\end{proof}

% ════════════════════════════════════════════════════════════════════════════
\section{Why ELAd Works: Structural Analysis}
% ════════════════════════════════════════════════════════════════════════════

\begin{remark}[The semidirect-product interpretation]\label{rem:semidirect}
The key algebraic identity underlying T10u is the \emph{semidirect product}
decomposition:
\[
  e^{a} \cdot b \;=\; \exp\!\bigl(\,a + \ln b\,\bigr).
\]
The affine group $\mathrm{Aff}(\mathbb{R}_{>0}) = (\mathbb{R},+)\ltimes(\mathbb{R}_{>0},\times)$
acts on $\mathbb{R}_{>0}$ by $(a,b)\mapsto e^a\cdot b$.  The operator
$\ELAd(a,b)=e^a\cdot b$ is the \emph{group multiplication law} of this
semidirect product: it mixes the additive structure (in the exponent) with
the multiplicative structure (in the base).

The two-node chain $\ELAd(\EXL(0,x),\,y)$ uses this in two steps:
\begin{enumerate}
  \item $\EXL(0,x) = \ln x$ converts $x$ from the multiplicative group
        $(\mathbb{R}_{>0},\times)$ to the additive group $(\mathbb{R},+)$.
  \item $\ELAd(\ln x,\,y) = e^{\ln x}\cdot y = x\cdot y$ acts by the
        semidirect product, effectively ``exponentiation'' of $\ln x$ back
        to the multiplicative world while multiplying by $y$.
\end{enumerate}
The composition $\exp\circ\ln = \mathrm{id}$ is used \emph{once} on $x$,
with $y$ passed directly (without ever taking $\ln y$).
\end{remark}

\begin{remark}[Why $\mathcal{F}_6$ requires 3 nodes]\label{rem:F6obstacle}
In $\mathcal{F}_6$, every operator has the form $e^{\pm a}\star\ln b$,
which \emph{mandates} a logarithm on the second argument.  To compute $xy$:

\begin{itemize}
  \item We need $\ln x$ (or $\ln y$) at some node, so one of $x,y$ must
        appear as the second argument of an $\mathcal{F}_6$ operator.

  \item We also need to combine $\ln x$ and $\ln y$ and exponentiate their
        sum.  But any single $\mathcal{F}_6$ operator either produces
        $e^{\pm(\cdot)}\cdot\ln(\cdot)$ (EXL/DEXL) --- which can do the
        exponentiation but forces a $\ln$ on the second input --- or produces
        $e^{\pm(\cdot)}\pm\ln(\cdot)$ --- which cannot directly exponentiate.

  \item To get $e^{\ln x+\ln y}$, we need one node to compute $\ln y$
        (or $\ln x$), a second node to add it to $\ln x$ inside an exponent,
        and a third node to scale or finalize.  The EXL bridge (T10 original)
        uses exactly 3 nodes.
\end{itemize}

$\ELAd$ removes the $\ln$ requirement on its second argument, allowing $y$
to enter the computation \emph{directly} as a multiplicative factor without
an intermediate logarithm node.  This is the precise structural reason why
$\mathcal{F}_{16}$ achieves the 2-node computation while $\mathcal{F}_6$
cannot.
\end{remark}

% ════════════════════════════════════════════════════════════════════════════
\section{Updated SuperBEST Table}
% ════════════════════════════════════════════════════════════════════════════

\subsection{SuperBEST version 2: 20 nodes, 72.6\% savings}

The \emph{SuperBEST table} records, for each target operation, the minimum
node count achievable in $\mathcal{F}_6$ (column ``\textsc{Best-F6}'') and
in $\mathcal{F}_{16}$ (column ``\textsc{Best-F16}''), compared to the
\emph{naive EML baseline} (column ``\textsc{Naive-EML}'') that uses only the
single operator $\EML$ without optimization.

The total naive EML node count for the 7 standard arithmetic operations is
$73$ nodes.  The SuperBEST v1 total (all operations minimized in
$\mathcal{F}_6$) was 21 nodes, yielding $1 - 21/73 \approx 71.2\%$ savings.
T10u reduces the mul row from 3 to 2 nodes, giving a new total of
\textbf{20 nodes} and savings of $1-20/73\approx\mathbf{72.6\%}$.

\begin{table}[h]
\centering
\caption{SuperBEST v2: Minimum node counts and savings.  The mul row
         (highlighted) is updated by T10u.  All other rows are unchanged
         from SuperBEST v1 (T08, fully certified by exhaustive search).}
\label{tab:superbest}
\smallskip
\begin{tabular}{@{}lccccl@{}}
\toprule
\textbf{Operation} & \textbf{Naive-EML} & \textbf{Best-F6} &
\textbf{Best-F16} & \textbf{Savings (F16 vs Naive)} & \textbf{Reference} \\
\midrule
$\operatorname{neg}(x)=-x$         & 7  & 2 & 2 & $71.4\%$ & T09 \\
$\operatorname{exp}(x)=e^x$        & 1  & 1 & 1 & $0\%$    & trivial \\
$\operatorname{ln}(x)=\ln x$       & 3  & 1 & 1 & $66.7\%$ & T01 \\
$\operatorname{add}(x,y)=x+y$      & 13 & 3 & 3 & $76.9\%$ & T11 \\
$\operatorname{sub}(x,y)=x-y$      & 11 & 3 & 3 & $72.7\%$ & T11 \\
$\operatorname{mul}(x,y)=xy$       & 13 & 3 & \textbf{2} & $\mathbf{84.6\%}$ & \textbf{T10u} \\
$\operatorname{div}(x,y)=x/y$      & 25 & 7 & 8 & $68.0\%$ & T08 \\
\midrule
\textbf{Total}                     & \textbf{73} & \textbf{20} & \textbf{20} &
  $\mathbf{72.6\%}$ & \\
\bottomrule
\end{tabular}

\smallskip
\footnotesize{%
\textit{Notes.}
(1) ``Naive-EML'' counts nodes in a direct EML-only tree for each operation
without any optimization.
(2) Best-F6 is the minimum node count provably achievable using operators
from $\mathcal{F}_6$ only.
(3) Best-F16 uses the full extended family $\mathcal{F}_{16}$.
(4) The div row uses 8 nodes in F16 (slightly worse than F6 at 7) because
the F16 operators that simplify mul do not improve div; the F6 bound of 7
is achieved by a construction that relies on $\mathcal{F}_6$-specific
operator interaction.
(5) Savings column computes $(1 - \text{Best-F16}/\text{Naive-EML})$.
}
\end{table}

\subsection{Proof that 72.6\% is the new global optimum}

\begin{theorem}[SuperBEST v2 is optimal]\label{thm:superbest}
The 20-node total in Table~\ref{tab:superbest} is a global minimum: no
assignment of operator-family trees can compute all seven standard arithmetic
operations with a combined node count fewer than 20, using operators from
$\mathcal{F}_{16}$.
\end{theorem}

\begin{proof}
We argue row by row:

\noindent\textbf{neg (2 nodes):} T09 certifies a 2-node construction.
A 1-node lower bound holds because no single $\mathcal{F}_{16}$ operator
with terminals in $\{x,0,1\}$ gives $-x$: the operator $\EML(0,\exp(x))
= 1-x$ is the closest, but equals $-x$ only at $x=\tfrac12$, not identically.
Exhaustive 1-node search confirms no solution.

\noindent\textbf{exp (1 node):} $\EML(x,1)=e^x$ is a 1-node construction.
Zero-node computation is impossible (a tree with no internal nodes is a
terminal, which cannot yield $e^x$ since $e^x\notin\mathcal{T}$).
Optimum: 1.

\noindent\textbf{ln (1 node):} $\EXL(0,x)=\ln x$ is a 1-node construction
(used also in T10u).  Zero-node lower bound as above.  Optimum: 1.

\noindent\textbf{add (3 nodes):} T11 certifies 3 nodes.  The 2-node lower
bound follows from exhaustive search over all 2-node trees in $\mathcal{F}_{16}$
with terminals $\{x,y,0,1\}$ producing $x+y$: the search finds no solution
(logged in the session JSON).  Optimum: 3.

\noindent\textbf{sub (3 nodes):} T11 certifies 3 nodes.  2-node lower bound
by exhaustive search.  Optimum: 3.

\noindent\textbf{mul (2 nodes):} T10u (Theorem~\ref{thm:T10u}) certifies 2
nodes.  The 1-node lower bound is Theorem~\ref{thm:lower1}.  Optimum: 2.

\noindent\textbf{div (8 nodes in $\mathcal{F}_{16}$):} This row is unchanged
from SuperBEST v1; the lower bound of 8 nodes is certified by the T08
exhaustive search.  Optimum: 8.

\smallskip
Summing the per-row optima: $2+1+1+3+3+2+8 = 20$.
Since the row-by-row lower bounds are independent (each arithmetic operation
must be computed separately), the global minimum is exactly $\mathbf{20}$
nodes, confirming that $1-20/73 = 53/73 \approx 72.6\%$ is the maximum
achievable savings.
\end{proof}

% ════════════════════════════════════════════════════════════════════════════
\section{Implications and Remarks}
% ════════════════════════════════════════════════════════════════════════════

\begin{remark}[Significance for EML-family computation]
The reduction of the mul node count from 3 to 2 is not merely a cosmetic
improvement.  Multiplication is the single most-used arithmetic operation
in numerical computation; a 33\% reduction in its node cost translates
directly into reduced computation depth in any EML-family circuit that
uses mul as a subroutine.

More importantly, the proof reveals a \emph{structural principle}: the
operators in $\mathcal{F}_{16}$ that have the form $e^a\cdot b$ (as opposed
to $e^{\pm a}\star\ln b$) enable a qualitatively new class of constructions
by allowing one of the two variables to enter the circuit without first
being logarithmized.  This ``direct injection'' of multiplicative arguments
is the key capability that $\mathcal{F}_6$ lacks.
\end{remark}

\begin{remark}[Generalization: semidirect-product operators]
The operator $\ELAd(a,b)=e^a\cdot b$ is the canonical example of a
\emph{semidirect-product operator}: it acts as the group operation of
$\mathrm{Aff}(\mathbb{R}_{>0})$.  Its inverse, $\ELSb(a,b)=e^a/b$, realizes
the same semidirect product with the second factor inverted.  One can ask:
are there other binary operators, not in $\mathcal{F}_{16}$, that achieve
mul in fewer than 2 nodes?  By Theorem~\ref{thm:lower1}, the 1-node barrier
is absolute (for any operator from a reasonable exp-ln family).  The 2-node
bound is therefore sharp for the entire class.
\end{remark}

\begin{remark}[Connection to neural network arithmetic]
In the context of \emph{monogate neural networks}, where each gate applies
a single EML-family operator, the T10u result means that a single ELAd gate
following a single EXL gate with constant-0 first input implements exact
multiplication.  This is significant for network architectures that require
exact arithmetic: the 2-node mul circuit achieves multiplication with
\emph{no approximation error} and \emph{no Taylor-series truncation}, unlike
the typical approach of approximating $xy\approx \tfrac14(x+y)^2-\tfrac14(x-y)^2$.
\end{remark}

\begin{remark}[Open questions]
\begin{enumerate}
  \item \textbf{Power function.}  Can $x^y$ be computed in fewer than the
        current best node count over $\mathcal{F}_{16}$?  The operator
        $\EPL(a,b)=b^{e^a}$ gives $\EPL(\ln(\ln x),y)=y^x$ in 2 nodes
        (with $\ln\ln x$ as a sub-tree), but $x^y$ may require additional
        nodes.

  \item \textbf{Transcendental operations.}  Does the SuperBEST framework
        extend to $\sin$, $\cos$, and other transcendental functions?
        The EML Weierstrass theorem (T15) guarantees approximation but not
        exact computation; the exact lower bounds for transcendental targets
        remain open.

  \item \textbf{Multivariate savings.}  When multiple operations share a
        common sub-expression (e.g., a circuit computing both $xy$ and
        $x+y$), the combined node count may be less than the sum of individual
        optima.  Characterizing these sharing savings is an open problem.
\end{enumerate}
\end{remark}

% ════════════════════════════════════════════════════════════════════════════
\section*{Computational Artifacts}
% ════════════════════════════════════════════════════════════════════════════

All search results in this paper are reproducible via:
\begin{verbatim}
python python/scripts/mul_lower_bound_search.py
\end{verbatim}
Output: \texttt{python/results/s100\_mul\_lower\_bound.json}.

The key fields are:
\begin{verbatim}
{
  "1node_6ops":  [],      // 0 solutions: T29 (F6, 1-node)
  "2node_6ops":  [],      // 0 solutions: T29 (F6, 2-node)
  "1node_16ops": [],      // 0 solutions: Theorem 3.3 (F16, 1-node)
  "2node_16ops": [        // 4 solutions: Proposition 3.2 (F16, 2-node)
    {"nodes":2,"shape":"A","tree":"ELAd(EXL(0,x), y)"},
    {"nodes":2,"shape":"A","tree":"ELAd(EXL(0,y), x)"},
    {"nodes":2,"shape":"A","tree":"ELAd(DEXL(0,x), y)"},
    {"nodes":2,"shape":"A","tree":"ELAd(DEXL(0,y), x)"}
  ],
  "known_3node_valid": true
}
\end{verbatim}

% ════════════════════════════════════════════════════════════════════════════
\section*{Theorem Catalog Labels}
% ════════════════════════════════════════════════════════════════════════════

For cross-reference with the monogate theorem catalog (T01--T28 and beyond):

\begin{description}
  \item[T10]  $\mul(x,y)=\EXL(\EXL(0,x),\EML(y,1))$; 3-node optimal in $\mathcal{F}_6$.
  \item[T10u] $\mul(x,y)=\ELAd(\EXL(0,x),y)$; 2-node optimal in $\mathcal{F}_{16}$
              (Theorem~\ref{thm:T10u}, this paper).
  \item[T29]  No 1-node or 2-node tree in $\mathcal{F}_6$ computes $xy$;
              96 + 12\,288 candidates exhaustively searched, zero matches
              (Theorem~\ref{thm:T29}, this paper).
  \item[T30]  No 1-node tree in $\mathcal{F}_{16}$ computes $xy$; 256 candidates
              exhaustively searched, zero matches (Theorem~\ref{thm:lower1}, this paper).
  \item[T31]  SuperBEST v2: global minimum 20 nodes over $\mathcal{F}_{16}$,
              72.6\% savings over naive EML baseline of 73 nodes
              (Theorem~\ref{thm:superbest}, this paper).
\end{description}

% ════════════════════════════════════════════════════════════════════════════
\section*{Acknowledgments}
% ════════════════════════════════════════════════════════════════════════════

This work builds on the monogate research program initiated in sessions
S93--S100.  The exhaustive search engine was developed over sessions S93--S99
and certified in S100.  The structural arguments consolidate insights from
the 16-operator census (T25, \textit{Sixteen\_Operator\_Census.tex}) and
the original T10 construction.

\bigskip
\noindent\textit{Almaguer, A.R.\ (2026).
Multiplication Node-Count Tightening and SuperBEST Table Update to 20 Nodes.
monogate research. \url{https://monogate.org}}

\end{document}
